\chapter[Band 14: Paarinfima und Paarsuprema]{Band 14 auf einen Blick}
\noindent\textbf{Paarinfima und Paarsuprema}\par\medskip
\noindent\emph{Lesart.} Es gilt \(\PartOrd{A,\leq}\).
Zur knappen Darstellung schreiben wir in dieser Übersicht
\(s(x,y)=\sup_{A,\leq}(\{x,y\})\) und
\(i(x,y)=\inf_{A,\leq}(\{x,y\})\), jeweils nur wenn der
betreffende Wert existiert. Eine partielle Ordnung garantiert
seine Existenz im Allgemeinen noch nicht.
\begin{BandUebersicht}
\textbf{Existenz eines Paarsupremums}
\(\operatorname{SupEx}_{A,\leq}(x,y)\) bedeutet:
\(x,y\in A\), und die oberen Schranken von \(\{x,y\}\)
besitzen ein kleinstes Element.
\UebFormel{\begin{aligned}
&\exists s\in\UB_{A,\leq}(\{x,y\})\\
&\quad\forall u\in\UB_{A,\leq}(\{x,y\})\,(s\leq u).
\end{aligned}}
\UebQuellen{\FormulaRefAuto{PairSupremumExistence}[def]}
&
\textit{Universelle Eigenschaft}
Unter dieser Existenzvoraussetzung:
\UebFormel{\begin{aligned}
&x\leq s(x,y),\qquad y\leq s(x,y),\\
&u\in A\dsep x\leq u\dsep y\leq u\vdash s(x,y)\leq u.
\end{aligned}}
\UebQuellen{\FormulaRefAuto{PairSupremumFirstOperand}[thm];
\FormulaRefAuto{PairSupremumSecondOperand}[thm];
\FormulaRefAuto{PairSupremumUniversal}[thm]}\\
\addlinespace[.8ex]
\textbf{Existenz eines Paarinfimums}
\(\operatorname{InfEx}_{A,\leq}(x,y)\) bedeutet:
\(x,y\in A\), und die unteren Schranken von \(\{x,y\}\)
besitzen ein größtes Element.
\UebFormel{\begin{aligned}
&\exists i\in\LB_{A,\leq}(\{x,y\})\\
&\quad\forall d\in\LB_{A,\leq}(\{x,y\})\,(d\leq i).
\end{aligned}}
\UebQuellen{\FormulaRefAuto{PairInfimumExistence}[def]}
&
\textit{Duale universelle Eigenschaft}
Unter dieser Existenzvoraussetzung:
\UebFormel{\begin{aligned}
&i(x,y)\leq x,\qquad i(x,y)\leq y,\\
&d\in A\dsep d\leq x\dsep d\leq y\vdash d\leq i(x,y).
\end{aligned}}
\UebQuellen{\FormulaRefAuto{PairInfimumFirstOperand}[thm];
\FormulaRefAuto{PairInfimumSecondOperand}[thm];
\FormulaRefAuto{PairInfimumUniversal}[thm]}\\
\addlinespace[.8ex]
\textbf{Wiederholte und vertauschte Operanden}
Die Paarmengen \(\{x,x\}\) und \(\{x,y\}=\{y,x\}\)
verbinden Mengengesetze mit den Extremaldefinitionen.
&
\textit{Idempotenz und Kommutativität}
Für \(x\in A\) existieren die wiederholten Paarwerte:
\UebFormel{s(x,x)=x,\qquad i(x,x)=x.}
Für existierende Werte gilt
\UebFormel{s(x,y)=s(y,x),\qquad i(x,y)=i(y,x).}
\UebQuellen{\FormulaRefAuto{PairSupremumIdempotent}[thm];
\FormulaRefAuto{PairInfimumIdempotent}[thm];
\FormulaRefAuto{PairSupremumCommutative}[thm];
\FormulaRefAuto{PairInfimumCommutative}[thm]}\\
\addlinespace[.8ex]
\textbf{Geschachtelte Paarwerte}
Vorausgesetzt werden für die Supremumsform alle vier Existenzen
\(s(x,y)\), \(s(y,z)\), \(s(s(x,y),z)\), \(s(x,s(y,z))\).
Für die Infimumsform werden entsprechend alle vier Infima vorausgesetzt.
&
\textit{Assoziativität unter den Existenzbedingungen}
\UebFormel{\begin{aligned}
&s(s(x,y),z)=s(x,s(y,z)),\\
&i(i(x,y),z)=i(x,i(y,z)).
\end{aligned}}
\UebQuellen{\FormulaRefAuto{PairSupremumAssociative}[thm];
\FormulaRefAuto{PairInfimumAssociative}[thm]}\\
\addlinespace[.8ex]
\textbf{Inklusionsordnung auf einer Mengenfamilie}
Sei \(M\subseteq\powerset(U)\). Die Relation ist die
gewöhnliche Teilmengenrelation; für \(X,Y\in M\) sind
\(X\cap Y\) und \(X\cup Y\) zunächst nur Teilmengen von \(U\).
&
\textit{Schnitt und Vereinigung als Paarwerte}
\UebFormel{\begin{aligned}
&\PartOrd{M,\subseteq},\\
&X\cap Y\in M\vdash\inf_{M,\subseteq}(\{X,Y\})=X\cap Y,\\
&X\cup Y\in M\vdash\sup_{M,\subseteq}(\{X,Y\})=X\cup Y.
\end{aligned}}
\UebQuellen{\FormulaRefAuto{InclusionOrderOnFamily}[thm];
\FormulaRefAuto{InclusionFamilyPairInfimum}[thm];
\FormulaRefAuto{InclusionFamilyPairSupremum}[thm]}\\
\addlinespace[.8ex]
\textbf{Potenzmenge als abgeschlossener Träger}
Für \(M=\powerset(U)\) sind Schnitt und Vereinigung stets
wieder Elemente des Trägers.
&
\textit{Paarwerte in der Potenzmenge}
Für \(X,Y\subseteq U\):
\UebFormel{\begin{aligned}
&\inf_{\powerset(U),\subseteq}(\{X,Y\})=X\cap Y,\\
&\sup_{\powerset(U),\subseteq}(\{X,Y\})=X\cup Y.
\end{aligned}}
\UebQuellen{\FormulaRefAuto{PowerSetPairInfimumIntersection}[thm];
\FormulaRefAuto{PowerSetPairSupremumUnion}[thm]}\\
\end{BandUebersicht}
